\documentclass[12pt,a4paper]{article}

\usepackage[utf8]{inputenc}
\usepackage[T1]{fontenc}
\usepackage[italian]{babel}
\usepackage{amsmath,amssymb,amsthm}
\usepackage{geometry}
\usepackage{booktabs}
\usepackage{array}
\usepackage{xcolor}
\usepackage{hyperref}
\usepackage[expansion=false]{microtype}
\usepackage{parskip}
\usepackage{enumitem}
\usepackage{tcolorbox}
\usepackage{mathtools}

\geometry{margin=2.5cm, top=3cm, bottom=3cm}

\hypersetup{
  colorlinks=true,
  linkcolor=blue!70!black,
  urlcolor=blue!70!black,
}

\tcbuselibrary{skins, breakable}

% Box per definizioni
\newtcolorbox{definizione}[1][]{
  colback=blue!5!white, colframe=blue!60!black,
  fonttitle=\bfseries, title=Definizione,
  breakable, #1
}

% Box per osservazioni importanti
\newtcolorbox{osservazione}[1][]{
  colback=orange!8!white, colframe=orange!70!black,
  fonttitle=\bfseries, title=Osservazione,
  breakable, #1
}

% Box per il riepilogo
\newtcolorbox{riepilogo}[1][]{
  colback=green!5!white, colframe=green!60!black,
  fonttitle=\bfseries, title=Riepilogo,
  breakable, #1
}

% Shorthand
\newcommand{\Rt}{R(t)}
\newcommand{\lbase}{\lambda_{\mathrm{base}}}
\newcommand{\loglik}{\ell}
\newcommand{\Smax}{S_{\max}}

\title{%
  \textbf{Stima di $R(t)$ con il metodo \mbox{Monte Carlo -- Metropolis}}\\[0.5em]
  \large Materiale di studio generato per approfondire i fondamenti teorici\\[0.3em]
  \normalsize Esame MSN --- Modellazione e Simulazioni Numeriche
}
\author{Samuel Seligardi \& Claude Code}
\date{A.A. 2025/2026}

\begin{document}

\maketitle
\tableofcontents
\newpage

% ===========================================================================
\section{Contesto e motivazione}
% ===========================================================================

Il \textbf{numero di riproduzione effettivo} $R(t)$ è la quantità epidemiologica centrale
per monitorare un'epidemia: rappresenta il numero medio di nuovi infetti generati da un
singolo caso infettivo nel giorno $t$, tenendo conto delle misure di contenimento e
dell'immunità parziale della popolazione.

\begin{itemize}
  \item $R(t) > 1$: l'epidemia è in crescita (ogni caso genera in media più di un secondario)
  \item $R(t) = 1$: l'epidemia è stazionaria
  \item $R(t) < 1$: l'epidemia è in regressione
\end{itemize}

Il Metodo~1 (SIR) stima $R(t)$ tramite la formula approssimata
$R(t) \approx 1 + T_R \cdot \Delta I / I$, che dipende dal parametro $T_R = 1/\gamma$
e introduce un ritardo strutturale di circa una settimana.

Il \textbf{Metodo~2} (implementato dalla metodologia ISS/FBK, Merler 2020)
usa un approccio completamente diverso:
\begin{enumerate}
  \item modella esplicitamente la \emph{distribuzione dei tempi di generazione} $\phi(s)$;
  \item usa un modello statistico Poisson per i nuovi contagi;
  \item stima $R(t)$ tramite un algoritmo Metropolis applicato \emph{indipendentemente}
        per ogni giorno.
\end{enumerate}

Vantaggi principali rispetto al Metodo~1:
\begin{itemize}
  \item non richiede la conoscenza di $\gamma$ (o $T_R$);
  \item fornisce una stima che anticipa il dato INFN di circa una settimana;
  \item è più accurato (errore massimo $0.20$--$0.27$ vs $0.2$--$0.6$ del SIR).
\end{itemize}

% ===========================================================================
\section{Il numero di riproduzione: dall'equazione di rinnovo}
% ===========================================================================

\subsection{L'equazione di Lotka--Euler}

Il punto di partenza del Metodo~2 è l'equazione di rinnovo (nota anche come
\emph{equazione di Lotka--Euler} in demografia):
\begin{equation}
  C(t) = R(t) \sum_{s=1}^{\infty} \phi(s)\, C(t-s)
  \label{eq:rinnovo}
\end{equation}
dove:
\begin{itemize}
  \item $C(t)$ = nuovi casi (sintomatici) al giorno $t$;
  \item $\phi(s)$ = probabilità che un contagio avvenga $s$ giorni dopo l'infezione
        del caso primario (\emph{distribuzione del tempo di generazione});
  \item $R(t)$ = numero di riproduzione al giorno $t$, la quantità da stimare.
\end{itemize}

\begin{osservazione}
L'equazione~\eqref{eq:rinnovo} non è un'approssimazione: è la definizione operativa
di $R(t)$ nel contesto dei modelli a rinnovo (renewal models). Il modello SIR ne è
un caso speciale con $\phi(s) = \gamma e^{-\gamma s}$ (distribuzione esponenziale
con parametro $\gamma$).
\end{osservazione}

\subsection{Riscrittura con la base}

Per un giorno $t$ fissato, il termine $\sum_s \phi(s)\,C(t-s)$ è una costante nota
(dipende solo dalla storia passata, non da $R(t)$). Lo chiamiamo \textbf{base}:
\begin{equation}
  \lbase(t) \;=\; \sum_{s=1}^{\Smax} \phi(s)\, C(t-s)
  \label{eq:lbase}
\end{equation}
dove il troncamento a $\Smax = 25$ giorni è giustificato perché $\phi(s) < 2 \times 10^{-3}$
per $s > 25$ (si veda la Sezione~\ref{sec:phi}).

L'equazione di rinnovo~\eqref{eq:rinnovo} diventa allora:
\begin{equation}
  \boxed{\lambda(t) \;=\; R(t) \cdot \lbase(t)}
  \label{eq:lambda}
\end{equation}
dove $\lambda(t)$ è il \emph{numero atteso} di nuovi casi al giorno $t$.

% ===========================================================================
\section{La distribuzione dei tempi di generazione $\phi(s)$}
\label{sec:phi}
% ===========================================================================

\subsection{Significato biologico}

Il \textbf{tempo di generazione} $X$ è la variabile aleatoria che misura il ritardo,
in giorni, tra il momento in cui un individuo viene infettato e il momento in cui
infetta a sua volta un caso secondario. Non è direttamente osservabile: si stima
indirettamente dall'analisi di cluster di casi con coppie infettante--infettato tracciate.

La distribuzione $\phi(s) = P(X = s)$ (discreta, definita su $s = 1, 2, 3, \ldots$)
è cruciale perché determina quanta ``memoria storica'' entra nel calcolo di $\lambda(t)$:
un $\phi$ con coda lunga significa che casi molto vecchi contribuiscono ancora
alla trasmissione odierna.

\subsection{La distribuzione Gamma continua}

Guzzetta e Merler (7 dicembre 2020) stimano per COVID-19 in Italia che il tempo di
generazione \emph{continuo} $X$ segue una distribuzione Gamma:
\begin{equation}
  X \sim \mathrm{Gamma}(\alpha,\, \theta) \qquad \alpha = 1.87, \quad \theta = \frac{1}{\beta} = \frac{1}{0.28} \approx 3.57\ \text{giorni}
\end{equation}
con densità di probabilità:
\begin{equation}
  f(x;\,\alpha,\theta) \;=\; \frac{x^{\alpha-1}\, e^{-x/\theta}}{\Gamma(\alpha)\, \theta^{\alpha}}, \qquad x > 0
  \label{eq:gamma_pdf}
\end{equation}
dove $\Gamma(\alpha) = \int_0^\infty t^{\alpha-1} e^{-t}\, dt$ è la \textbf{funzione Gamma
di Eulero} (non la distribuzione Gamma: attenzione alla notazione sovraccarica).

I momenti della distribuzione Gamma sono:
\begin{align}
  \mathbb{E}[X] &= \alpha \theta \approx 1.87 \times 3.57 \approx 6.68\ \text{giorni} \\
  \mathrm{Var}[X] &= \alpha \theta^2 \approx 1.87 \times 12.74 \approx 23.8\ \text{giorni}^2 \\
  \sigma_X &= \sqrt{\mathrm{Var}[X]} \approx 4.88\ \text{giorni}
\end{align}

La distribuzione è \textbf{asimmetrica a destra} (il coefficiente di asimmetria vale
$2/\sqrt{\alpha} \approx 1.46$): la maggior parte dei contagi avviene entro 5--10 giorni, con una coda
che si estende fino a 20--25 giorni. Questo è tipico dei processi biologici: esistono
casi rari di trasmissione tardiva (contatti prolungati, casi gravi con lunga infettività).

\begin{osservazione}[title=Perché la Gamma?]
La distribuzione Gamma è la scelta naturale per modellare tempi positivi con asimmetria
a destra. È inoltre la somma di $\alpha$ variabili esponenziali indipendenti con parametro
$\beta$ (per $\alpha$ intero), il che le conferisce una interpretazione biologica:
il processo di trasmissione può avvenire in $\alpha$ ``stadi'' successivi.
Per $\alpha = 1$ si riduce all'esponenziale (il caso del modello SIR compartimentale).
\end{osservazione}

\subsection{Discretizzazione}

Il modello lavora con \textbf{giorni interi}, ma $X$ è continua. La probabilità discreta
$\phi(s)$ si ottiene integrando $f$ sull'intervallo unitario centrato in $s$:
\begin{equation}
  \phi(s) \;=\; P\!\left(s - \tfrac{1}{2} \le X \le s + \tfrac{1}{2}\right)
          \;=\; \int_{s-1/2}^{s+1/2} f(u;\,\alpha,\theta)\, du, \qquad s = 1, 2, \ldots
  \label{eq:phi_discreta}
\end{equation}

Questa è la discretizzazione di \emph{midpoint binning}: ogni valore intero cattura
l'intera probabilità dell'intervallo unitario centrato su di esso. È il metodo corretto
per discretizzare una densità continua (a differenza di campionare $f(s)$ puntualmente,
che non dà una distribuzione di probabilità).

\textbf{Troncamento:} per $s > 25$, il contributo $\phi(s)$ è trascurabile. Si verifica
che $\sum_{s=26}^\infty \phi(s) \approx 5.2 \times 10^{-3}$, cioè lo 0.52\% della
probabilità totale. Si pone quindi $\Smax = 25$ e si \textbf{rinormalizza}:
\begin{equation}
  \phi(s) \;\leftarrow\; \frac{\phi(s)}{\sum_{k=1}^{25} \phi(k)}, \qquad s = 1, \ldots, 25
\end{equation}
affinché il vettore $(\phi(1), \ldots, \phi(25))$ sia una distribuzione di probabilità discreta.

\subsection{Valori numerici attesi}

\begin{center}
\begin{tabular}{cccccccccc}
\toprule
$s$ (giorni) & 1 & 2 & 3 & 4 & 5 & 6 & 7 & 8 & $\cdots$ \\
\midrule
$\phi(s)$ (approx.) & 0.073 & 0.103 & 0.111 & 0.108 & 0.099 & 0.088 & 0.076 & 0.064 & $\cdots$ \\
\bottomrule
\end{tabular}
\end{center}

La media $\sum_s s\,\phi(s) \approx 6.65$ giorni deve coincidere con la media teorica
della Gamma (a meno dell'errore di troncamento). Questa è la prima verifica di
correttezza da eseguire in MATLAB.

% ===========================================================================
\section{Il modello statistico: distribuzione di Poisson}
% ===========================================================================

\subsection{Motivazione}

Dato $R(t)$ e la storia passata $C(t-1), \ldots, C(t-25)$, il numero di nuovi casi
$C(t)$ è una variabile aleatoria. Il modello adottato è:
\begin{equation}
  C(t) \mid R(t) \;\sim\; \mathrm{Poisson}\!\left(\lambda(t)\right) = \mathrm{Poisson}\!\left(R(t)\cdot\lbase(t)\right)
  \label{eq:poisson}
\end{equation}

La distribuzione di Poisson è giustificata da due argomenti:

\begin{enumerate}
  \item \textbf{Argomento del processo puntuale}: in ogni istante, un individuo infettivo
        ha una piccola probabilità di trasmettere il virus. Se i contatti sono indipendenti
        e la probabilità di trasmissione è piccola ma il numero di individui suscettibili
        è grande, il numero di trasmissioni segue una Poisson (\emph{legge dei grandi numeri
        per processi di Bernoulli rari}).

  \item \textbf{Argomento della superspreading}: anche in presenza di eterogeneità
        (alcuni individui trasmettono molto di più di altri), il conteggio aggregato
        giornaliero si approssima bene con una Poisson per grandi popolazioni.
\end{enumerate}

La probabilità di massa della Poisson è:
\begin{equation}
  P\!\left(C(t) = c \;\middle|\; R(t)\right) \;=\; \frac{e^{-\lambda(t)}\, \lambda(t)^c}{c!}, \qquad c = 0, 1, 2, \ldots
\end{equation}

\begin{osservazione}
La Poisson ha un solo parametro $\lambda$ che funge sia da media che da varianza:
$\mathbb{E}[C] = \mathrm{Var}[C] = \lambda$. Ne consegue che la deviazione standard
relativa è $\sigma_C / \mathbb{E}[C] = 1/\sqrt{\lambda}$: la stima di $R(t)$ sarà
tanto più precisa quanto più grande è $\lambda(t)$, ossia quanto più casi ci sono.
\end{osservazione}

% ===========================================================================
\section{La funzione di verosimiglianza e il massimo di verosimiglianza}
\label{sec:lik}
% ===========================================================================

\subsection{Log-verosimiglianza}

Per un giorno $t$ fissato, con $\lbase(t)$ nota, vogliamo stimare $R(t)$ dall'osservazione
di $C(t) = c$. La \textbf{funzione di verosimiglianza} è:
\begin{equation}
  L(R) \;=\; P\!\left(C(t) = c \;\middle|\; R\right) \;=\; \frac{e^{-R\lbase}\,(R\lbase)^c}{c!}
\end{equation}
dove, per brevità, scriviamo $\lbase = \lbase(t)$.

La \textbf{log-verosimiglianza} è:
\begin{align}
  \ell(R) &= \log L(R) = -R\lbase + c\,\log(R\lbase) - \log(c!) \notag \\
           &= c\,\log R \;+\; c\,\log\lbase \;-\; R\lbase \;-\; \log(c!)
  \label{eq:loglik}
\end{align}

I termini $c\,\log\lbase$ e $-\log(c!)$ sono costanti rispetto a $R$: non influenzano
la posizione del massimo.

\subsection{Massimo di verosimiglianza (MLE)}

Deriviamo $\ell(R)$ rispetto a $R$ e uguagliamo a zero:
\begin{equation}
  \frac{d}{dR}\,\ell(R) \;=\; \frac{c}{R} - \lbase \;=\; 0
\end{equation}
da cui:
\begin{equation}
  \boxed{R_{\mathrm{MLE}} \;=\; \frac{c}{\lbase} \;=\; \frac{C(t)}{\lbase(t)}}
  \label{eq:mle}
\end{equation}

Per verificare che sia un massimo (e non un minimo), calcoliamo la derivata seconda:
\begin{equation}
  \frac{d^2}{dR^2}\,\ell(R) \;=\; -\frac{c}{R^2} \;<\; 0 \quad \text{per } R > 0,\, c > 0
\end{equation}
Poiché la derivata seconda è strettamente negativa, $\ell(R)$ è \textbf{strettamente concava}
su $R > 0$, e $R_{\mathrm{MLE}}$ è l'unico massimo globale.

\begin{osservazione}
Il risultato $R_{\mathrm{MLE}} = C(t)/\lbase(t)$ ha un'interpretazione intuitiva diretta:
$R(t)$ è il rapporto tra i casi osservati oggi e il numero di casi atteso se $R(t) = 1$.
\end{osservazione}

\subsection{Perché non usare direttamente l'MLE?}

L'MLE è la stima puntuale ottimale per $c$ grande (caso asintotico). Per $c$ piccolo,
la distribuzione di Poisson è fortemente asimmetrica, e la \emph{media della distribuzione
a posteriori} è preferibile al massimo per due motivi.

\medskip
\noindent\textit{Nota di notazione (anticipazione):} nella Sezione~\ref{sec:bayes} verrà
introdotta formalmente la \textbf{distribuzione a posteriori} $\pi(R \mid c)$, ovvero
la densità di probabilità per il parametro $R$ \emph{dopo} aver osservato il dato $c$.
Qui la lettera $\pi$ (con argomento) non è la costante $\pi \approx 3.14\ldots$, ma
la notazione statistica standard per una densità di probabilità; il simbolo
``$\mid c$'' si legge ``dato $c$'' (condizionato all'osservazione $c$).

\begin{enumerate}
  \item \textbf{Asimmetria}: quando $c$ è piccolo, la distribuzione posteriore
        $\pi(R \mid c) \propto e^{\ell(R)}$ è asimmetrica a destra. La media è spostata
        verso valori più alti rispetto all'MLE (il massimo). Usare la media riduce
        l'errore quadratico medio (\emph{bias-variance tradeoff}).

  \item \textbf{Quantificazione dell'incertezza}: il campionamento Metropolis produce
        naturalmente una distribuzione di campioni da $\pi(R \mid c)$, permettendo di
        calcolare intervalli di confidenza senza ipotesi aggiuntive.
\end{enumerate}

Per $c \to \infty$, la Poisson converge a una Gaussiana e la media coincide con il
massimo: $\mathbb{E}_{\pi}[R \mid c] \to R_{\mathrm{MLE}}$ (per il Teorema del Limite Centrale).

% ===========================================================================
\section{Approccio Bayesiano: la distribuzione posteriore}
\label{sec:bayes}
% ===========================================================================

\subsection{Notazione: $\pi(\cdot)$ come densità di probabilità}

\begin{osservazione}[title=Disambiguazione: $\pi(\cdot)$ non è la costante $3.14\ldots$]
In questo documento --- e in quasi tutta la letteratura statistica --- la lettera
$\pi$ \emph{con un argomento} (come $\pi_0(R)$ o $\pi(R \mid c)$) indica una
\textbf{funzione di densità di probabilità}. Questo uso è completamente distinto dalla
costante matematica $\pi \approx 3.14159\ldots$

Quando scriviamo $\pi(R \mid c)$ intendiamo: \emph{``la densità di probabilità della
variabile aleatoria $R$, condizionata all'osservazione $c$''}. L'argomento a sinistra
della barra verticale è la variabile; quello a destra è il dato su cui condizioniamo.
\end{osservazione}

\subsection{Il framework Bayesiano: tre ingredienti}

Nell'approccio \textbf{frequentista} usato finora per l'MLE, il parametro $R(t)$ è
trattato come una costante ignota: si cerca il valore che massimizza $L(R)$.

Nell'approccio \textbf{Bayesiano}, invece, il parametro ignoto $R(t)$ è trattato come
una \emph{variabile aleatoria}: se ne descrive la distribuzione di probabilità, che viene
aggiornata man mano che si osservano i dati. I tre ingredienti sono:

\begin{enumerate}
  \item \textbf{Distribuzione a priori} (\emph{prior}): $\pi_0(R)$.\\
        È la densità di probabilità per $R(t)$ \emph{prima} di osservare il dato
        $C(t) = c$. Codifica la conoscenza (o l'ignoranza) iniziale sul parametro.
        Il pedice $0$ distingue il prior dalla distribuzione aggiornata.

  \item \textbf{Verosimiglianza}: $L(R) = P\!\left(C(t) = c \mid R\right)$.\\
        Già calcolata nella Sezione~\ref{sec:lik}: misura la compatibilità tra il
        valore del parametro $R$ e il dato osservato $c$. È una funzione di $R$,
        \emph{non} una probabilità per $R$.

  \item \textbf{Distribuzione a posteriori} (\emph{posterior}): $\pi(R \mid c)$.\\
        È la densità di probabilità per $R(t)$ \emph{dopo} aver osservato $c$.
        Combina l'informazione del prior con quella contenuta nel dato.
        È l'oggetto che Metropolis campionerà.
\end{enumerate}

Il legame tra le tre componenti è il \textbf{Teorema di Bayes}:
\begin{equation}
  \pi(R \mid c)
  \;=\;
  \frac{L(R)\cdot\pi_0(R)}{\displaystyle\int_0^{\infty} L(R')\cdot\pi_0(R')\,dR'}
  \label{eq:bayes_formula}
\end{equation}
Il denominatore $Z = \int_0^{\infty} L(R')\,\pi_0(R')\,dR'$ è la
\textbf{costante di normalizzazione} (chiamata anche \emph{evidenza}): è un numero
reale positivo che dipende solo dal dato $c$ e garantisce che
$\int_0^\infty \pi(R \mid c)\,dR = 1$. Poiché $Z$ è una costante rispetto a $R$,
si può scrivere in forma proporzionale:
\begin{equation}
  \pi(R \mid c) \;\propto\; L(R)\cdot\pi_0(R)
  \label{eq:bayes_prop}
\end{equation}
Questa forma è sufficiente per l'algoritmo Metropolis, che non ha mai bisogno
di calcolare $Z$ esplicitamente (vedi Sezione~\ref{sec:metro}).

\subsection{Prior uniforme e prior improprio}

Adottiamo un \textbf{prior uniforme} su $R > 0$:
\begin{equation}
  \pi_0(R) \;\propto\; 1, \qquad R > 0
\end{equation}
Questo prior esprime la \emph{completa assenza di informazione a priori}: nessun
valore di $R(t)$ è considerato più plausibile degli altri prima di guardare i dati.

Il termine \textbf{improprio} (\emph{improper prior}) indica che questo prior non è
una vera distribuzione di probabilità nel senso classico, perché
$\int_0^{\infty} 1\,dR = +\infty$: non è normalizzabile a~1. I prior impropri sono
comunque accettati in statistica Bayesiana a condizione che la distribuzione a
posteriori risultante sia \emph{propria} (abbia integrale finito). Verifichiamo
questa condizione nel paragrafo successivo.

\subsection{La distribuzione posteriore esplicita}

Sostituendo $\pi_0(R) \propto 1$ nella~\eqref{eq:bayes_prop}:
\begin{equation}
  \pi(R \mid c) \;\propto\; L(R)\cdot 1 \;=\; L(R)
  \label{eq:posteriore}
\end{equation}
Con prior uniforme, la posteriore \emph{coincide con la verosimiglianza}
(a meno della costante di normalizzazione). Sostituendo la forma esplicita di $L(R)$:
\begin{equation}
  \pi(R \mid c) \;\propto\; R^c\, e^{-R\lbase}, \qquad R > 0
\end{equation}
Questo è esattamente il nucleo di una distribuzione \textbf{Gamma}:
$R \mid c \;\sim\; \mathrm{Gamma}(c+1,\; 1/\lbase)$.

L'integrale $\int_0^\infty R^c e^{-R\lbase}\,dR = \Gamma(c+1)/\lbase^{c+1}$ è finito
per ogni $c \ge 0$, confermando che la posteriore è propria (nonostante il prior
fosse improprio).

La \textbf{media posteriore} (stimatore ottimo sotto perdita quadratica, detto
\emph{stimatore di Bayes}) è:
\begin{equation}
  \mathbb{E}[R \mid c] \;=\; \frac{c+1}{\lbase}
\end{equation}
mentre la \textbf{moda} della posteriore (che con prior uniforme coincide con l'MLE)
è $c/\lbase$. La differenza $+1/\lbase$ è trascurabile per $c \gg 1$ ma significativa
per $c$ piccolo.

\begin{osservazione}
In linea di principio, per questa distribuzione posteriore (Gamma) la stima esatta
potrebbe essere calcolata analiticamente senza Metropolis. Il Metropolis è usato
nella metodologia ISS perché il modello reale include correlazioni tra giorni diversi
e prior non uniformi che rendono la distribuzione posteriore intractable analiticamente.
Nella nostra implementazione semplificata (stima giorno per giorno indipendente),
il Metropolis approssima correttamente la Gamma posteriore.
\end{osservazione}

% ===========================================================================
\section{L'algoritmo Metropolis}
\label{sec:metro}
% ===========================================================================

\subsection{Principio generale: catene di Markov Monte Carlo}

L'algoritmo \textbf{Metropolis} (Metropolis et al., 1953) è il capostipite dei metodi
MCMC (Markov Chain Monte Carlo). Il problema che risolve è:

\begin{definizione}
Sia $\mathcal{X}$ lo spazio dei valori ammissibili per il parametro di interesse
(nel nostro caso $\mathcal{X} = \mathbb{R}_{>0}$, i reali positivi, poiché $R > 0$).
Sia $\pi: \mathcal{X} \to \mathbb{R}_{\ge 0}$ una \emph{distribuzione target} ---
la distribuzione da cui vogliamo campionare --- di cui si conosce solo la forma
funzionale, non la costante di normalizzazione.

Il problema è costruire una \textbf{catena di Markov}
$\{X^{(0)}, X^{(1)}, \ldots, X^{(N)}\}$ su $\mathcal{X}$, ovvero una successione
di valori tali che $X^{(k)}$ dipende dal passato solo tramite $X^{(k-1)}$, la cui
\textbf{distribuzione stazionaria} sia $\pi$. Con ``distribuzione stazionaria''
si intende che, per $k$ sufficientemente grande, la distribuzione di $X^{(k)}$
converge a $\pi$ indipendentemente dal punto di partenza $X^{(0)}$.
\end{definizione}

Nel nostro contesto: la distribuzione target $\pi$ è la \textbf{distribuzione a
posteriori} $\pi(R \mid c)$ introdotta nella Sezione~\ref{sec:bayes}, ovvero:
\begin{equation}
  \pi(R) \;\propto\; e^{\ell(R)} \;=\; e^{\,c\log R\,-\,R\lbase} \;=\; R^c\,e^{-R\lbase}
\end{equation}
(scriviamo $\pi(R)$ per brevità, con $c$ e $\lbase$ fissi per il giorno $t$ corrente).
La costante di normalizzazione $\int_0^\infty R^c e^{-R\lbase}\,dR$ non viene mai
calcolata: Metropolis la evita interamente, come si vede dal passo~(b) qui sotto.

\subsection{L'algoritmo passo per passo}

\textbf{Inizializzazione:} $R^{(0)} = 1.0$ (soglia epidemica, valore neutro).

\medskip
\textbf{Passo $k$} ($k = 1, 2, \ldots, N_{\mathrm{iter}}$):

\begin{enumerate}[label=\textbf{\alph*)}]
  \item \textbf{Proposta.} Si campiona un candidato:
    \begin{equation}
      R^* = R^{(k-1)} + U(-\delta,\, +\delta), \qquad \delta = 1.5
      \label{eq:proposta}
    \end{equation}
    dove $U(-\delta, +\delta)$ è una variabile uniforme sull'intervallo $(-\delta, +\delta)$.
    La distribuzione di proposta è \textbf{simmetrica}: $q(R^* \mid R) = q(R \mid R^*)$,
    proprietà essenziale per Metropolis (non sarebbe necessaria per Metropolis--Hastings,
    che è la versione più generale).

  \item \textbf{Log-ratio.} Se $R^* > 0$, si calcola:
    \begin{align}
      \log\rho &= \ell(R^*) - \ell(R^{(k-1)}) \notag \\
               &= \bigl[c\log(R^*\lbase) - R^*\lbase\bigr]
                  - \bigl[c\log(R^{(k-1)}\lbase) - R^{(k-1)}\lbase\bigr] \notag \\
               &= c\,\log\!\frac{R^*}{R^{(k-1)}} - \lbase\!\left(R^* - R^{(k-1)}\right)
      \label{eq:logrho}
    \end{align}
    Il termine $c\,\log\lbase$ è presente sia in $\ell(R^*)$ che in $\ell(R^{(k-1)})$
    e si cancella esattamente. Questo significa che $\lbase$ entra nel log-ratio
    solo attraverso il termine lineare $-\lbase(R^* - R^{(k-1)})$.

  \item \textbf{Accettazione/rifiuto.} Si pone:
    \begin{equation}
      R^{(k)} =
      \begin{cases}
        R^{(k-1)}              & \text{se } R^* \le 0 \text{ (rifiuto forzato)} \\
        R^*                    & \text{se } \log\rho \ge 0 \text{ (accettazione certa)} \\
        R^* \text{ con prob.\ } e^{\log\rho} & \text{se } \log\rho < 0 \\
        R^{(k-1)}              & \text{altrimenti (rifiuto)}
      \end{cases}
      \label{eq:accettazione}
    \end{equation}
    In pratica: si genera $u \sim U(0,1)$ e si accetta se $\log u < \log\rho$
    (equivalente a $u < e^{\log\rho}$, ma numericamente più stabile per evitare
    overflow di $e^x$ per $\log\rho$ molto negativo).
\end{enumerate}

\textbf{Burn-in:} i primi $N_{\mathrm{burnin}} = 500$ valori sono scartati.
Durante il burn-in la catena ``si avvicina'' alla zona di alta probabilità a partire
dal punto iniziale $R=1$ (che potrebbe essere lontano dalla moda della posteriore).

\textbf{Stima finale:}
\begin{equation}
  \widehat{R}(t) \;=\; \frac{1}{N_{\mathrm{iter}} - N_{\mathrm{burnin}}} \sum_{k=N_{\mathrm{burnin}}+1}^{N_{\mathrm{iter}}} R^{(k)}
  \;=\; \frac{1}{4500} \sum_{k=501}^{5000} R^{(k)}
  \label{eq:stima}
\end{equation}

\subsection{Interpretazione fisica del log-ratio}

La formula~\eqref{eq:logrho} ha un'interpretazione intuitiva chiara. Separando i due termini:

\begin{center}
\renewcommand{\arraystretch}{1.5}
\begin{tabular}{lll}
\toprule
Termine & Segno tipico & Effetto sulla catena \\
\midrule
$c\,\log(R^*/R^{(k-1)})$ & $> 0$ se $R^* > R$ & Spinge verso $R$ più alto se ci sono casi \\
$-\lbase(R^* - R^{(k-1)})$ & $< 0$ se $R^* > R$ & Frena l'aumento proporzionalmente a $\lbase$ \\
\midrule
Bilancio & 0 & Equilibrio in $R^* = R_{\mathrm{MLE}} = c/\lbase$ \\
\bottomrule
\end{tabular}
\end{center}

In $R = R_{\mathrm{MLE}} = c/\lbase$, il log-ratio per qualsiasi spostamento $\epsilon$ è:
\begin{equation}
  \log\rho = c\,\log\!\frac{R_{\mathrm{MLE}}+\epsilon}{R_{\mathrm{MLE}}} - \lbase\,\epsilon
  \approx c\,\frac{\epsilon}{R_{\mathrm{MLE}}} - \lbase\,\epsilon
  = \epsilon\!\left(\frac{c}{R_{\mathrm{MLE}}} - \lbase\right) = 0
\end{equation}
confermando che l'MLE è il punto stazionario.

% ===========================================================================
\section{Garanzia teorica: bilancio dettagliato e convergenza}
% ===========================================================================

\subsection{Il problema: costruire $W$ data $\Pi$}

L'algoritmo Metropolis nasce come soluzione a un problema generale che condivide
la stessa logica dei processi di Markov discreti. Si vuole calcolare il
\textbf{valore di aspettazione} di una funzione $f$ rispetto a una distribuzione
target $\Pi$:
\begin{equation}
  \langle f(X) \rangle \;=\; \sum_i f(X_i)\,\Pi_i
  \label{eq:aspettazione_mc}
\end{equation}
dove la somma è su tutte le possibili configurazioni $X_i$, pesata dalla
distribuzione di probabilità $\Pi$.

Si costruisce a tal fine una \textbf{catena di Markov} con matrice di transizione $W$,
i cui elementi $W_{ij} = P(i \leftarrow j)$ sono probabilità di transizione dallo
stato $j$ allo stato $i$. Valgono le proprietà $W_{ij} \ge 0$ e
$\sum_i W_{ij} = 1$ per ogni $j$.

Iterando a partire da una distribuzione iniziale $P_0$, la distribuzione dopo $N$ passi è
$P^{(N)} = W^N P_0$. Se vale la condizione di \textbf{stazionarietà} $W\Pi = \Pi$
(ovvero $\Pi$ è l'autovettore di $W$ con autovalore $1$), allora per ergodicità:
\begin{equation}
  P^{(N)} = W^N P_0 \;\xrightarrow{\;N \to \infty\;}\; \Pi
\end{equation}
indipendentemente dal punto di partenza $P_0$. Grazie a questo, possiamo
\textbf{sostituire medie temporali} sull'evoluzione asintotica della catena con
\textbf{medie su $\Pi$}: questo è il principio fondamentale di tutti i metodi MCMC.

\begin{osservazione}[title=Problema fondamentale di MCMC]
\textbf{Data} la distribuzione target $\Pi$ (di cui si conosce la forma funzionale,
non la costante di normalizzazione), come si costruisce una matrice di transizione
$W$ tale che $W\Pi = \Pi$?
\end{osservazione}

Una \textbf{condizione sufficiente} (non necessaria) per soddisfare $W\Pi = \Pi$
è il bilancio dettagliato, derivato nel paragrafo seguente.

% ---------------------------------------------------------------------------
\subsection{Il bilancio dettagliato come condizione sufficiente}

La ragione per cui l'algoritmo Metropolis funziona è la soddisfazione della condizione
di \textbf{bilancio dettagliato} (detailed balance). Questa condizione garantisce che la
distribuzione target $\pi$ è una distribuzione stazionaria della catena di Markov.

\begin{definizione}[title=Definizione: bilancio dettagliato]
Una catena di Markov con probabilità di transizione $T(x \to y)$ soddisfa il bilancio
dettagliato rispetto a $\pi$ se:
\begin{equation}
  \pi(x)\, T(x \to y) \;=\; \pi(y)\, T(y \to x) \qquad \forall\, x, y \in \mathcal{X}
  \label{eq:bd}
\end{equation}
\end{definizione}

Per Metropolis, la probabilità di transizione (per proposta simmetrica) è:
\begin{equation}
  T(R \to R^*) \;=\; q(R^*\mid R)\cdot \alpha(R, R^*)
  \;=\; q(R^*\mid R)\cdot \min\!\left(1,\, \frac{\pi(R^*)}{\pi(R)}\right)
\end{equation}
dove i due fattori corrispondono alle due fasi sequenziali dell'algoritmo:
\begin{itemize}
  \item $q(R^*\mid R)$: \textbf{probabilità di proposta} --- con quale probabilità l'algoritmo
        estrae $R^*$ come candidato, dato che si trova in $R$. Nel nostro caso è la distribuzione
        uniforme $U(R-\delta,\, R+\delta)$: non dipende da $\pi$ e non sa nulla della
        distribuzione target.
  \item $\alpha(R, R^*)$: \textbf{probabilità di accettazione} --- dato che $R^*$ è stato proposto,
        con quale probabilità viene accettato. Questa dipende da $\pi$: vale $1$ se
        $R^*$ è più probabile del corrente, meno di $1$ altrimenti.
\end{itemize}
Il prodotto $q \cdot \alpha$ è la probabilità che entrambe le fasi abbiano successo,
ovvero che la catena si sposti effettivamente da $R$ a $R^*$.

\textbf{Verifica del bilancio dettagliato.} Poiché $q$ è simmetrica ($q(R^*\mid R) = q(R\mid R^*)$),
dividiamo per $q$ e verifichiamo:
\begin{equation}
  \pi(R)\cdot \min\!\left(1,\, \frac{\pi(R^*)}{\pi(R)}\right)
  \;=\;
  \pi(R^*)\cdot \min\!\left(1,\, \frac{\pi(R)}{\pi(R^*)}\right)
\end{equation}

\textbf{Caso 1}: $\pi(R^*) \ge \pi(R)$.
\begin{align*}
  \text{Lato sinistro} &= \pi(R)\cdot \frac{\pi(R^*)}{\pi(R)} = \pi(R^*) \\
  \text{Lato destro}  &= \pi(R^*)\cdot 1 = \pi(R^*) \quad \checkmark
\end{align*}

\textbf{Caso 2}: $\pi(R^*) < \pi(R)$ (simmetrico per scambio di $R$ e $R^*$). \checkmark

Il bilancio dettagliato implica che $\pi$ è una distribuzione stazionaria. Affinché
sia anche l'unica (condizione di \emph{ergodicità}), è sufficiente che la catena sia:
\begin{itemize}
  \item \textbf{irriducibile}: da qualsiasi punto si può raggiungere qualsiasi regione
        con probabilità positiva (garantito dalla proposta uniforme con $\delta = 1.5$,
        che copre un intervallo ampio intorno al valore corrente);
  \item \textbf{aperiodica}: non ci sono cicli deterministici (garantito dal fatto che
        ci sono proposte che vengono rifiutate, quindi la catena può rimanere ferma).
\end{itemize}

% ---------------------------------------------------------------------------
\subsection{Derivazione della probabilità di accettazione: $F(z) = \min(1,z)$}

Si fattorizza la probabilità di transizione in due contributi:
\begin{equation}
  W_{ij} \;=\; P^{(0)}_{ij} \cdot a_{ij}
\end{equation}
dove $P^{(0)}_{ij}$ è la \textbf{probabilità di proporre} la mossa $j \to i$
e $a_{ij} \in [0,1]$ è la \textbf{probabilità di accettare} la proposta.

Sostituendo nella condizione di bilancio dettagliato $W_{ij}\Pi_j = W_{ji}\Pi_i$:
\begin{equation}
  P^{(0)}_{ij}\, a_{ij}\, \Pi_j \;=\; P^{(0)}_{ji}\, a_{ji}\, \Pi_i
\end{equation}
da cui si ottiene il rapporto tra le probabilità di accettazione:
\begin{equation}
  \frac{a_{ij}}{a_{ji}} \;=\; \frac{P^{(0)}_{ji}\,\Pi_i}{P^{(0)}_{ij}\,\Pi_j}
  \label{eq:rapporto_acc}
\end{equation}

\textbf{Proposta simmetrica.} Se $P^{(0)}_{ij} = P^{(0)}_{ji}$ (come la distribuzione
$U(-\delta, +\delta)$ del nostro algoritmo), i termini di proposta si cancellano:
\begin{equation}
  \frac{a_{ij}}{a_{ji}} \;=\; \frac{\Pi_i}{\Pi_j}
\end{equation}

Si cerca dunque una funzione $F:[0,+\infty) \to [0,1]$ tale che, posto $z = \Pi_i/\Pi_j$:
\begin{equation}
  a_{ij} = F(z), \quad a_{ji} = F(z^{-1}), \qquad
  \text{con vincolo} \;\frac{F(z)}{F(z^{-1})} = z
  \label{eq:condF}
\end{equation}
La soluzione più semplice alla~\eqref{eq:condF} è:
\begin{equation}
  \boxed{F(z) \;=\; \min(1,\,z)}
  \label{eq:Fmin}
\end{equation}

\textbf{Verifica diretta:}
\begin{itemize}
  \item $z > 1$:\; $F(z) = 1$,\; $F(z^{-1}) = z^{-1}$;\; $F(z)/F(z^{-1}) = z$ \checkmark
  \item $z < 1$:\; $F(z) = z$,\; $F(z^{-1}) = 1$;\; $F(z)/F(z^{-1}) = z$ \checkmark
  \item $z = 1$:\; $F(z) = F(z^{-1}) = 1$;\; $F(z)/F(z^{-1}) = 1 = z$ \checkmark
\end{itemize}

L'algoritmo Metropolis risultante è dunque:
\begin{enumerate}
  \item Proponi la mossa $j \to i$ con distribuzione simmetrica
  \item Calcola $z = \Pi_i / \Pi_j$
  \item Accetta con probabilità $\min(1, z)$:
        \begin{itemize}
          \item se $\Pi_i \ge \Pi_j$: accetta sempre (stato proposto più probabile)
          \item se $\Pi_i < \Pi_j$: accetta con probabilità $\Pi_i/\Pi_j < 1$
        \end{itemize}
\end{enumerate}

Intuitivamente: le mosse verso stati più probabili sono sempre accettate; quelle verso
stati meno probabili sono accettate tanto meno quanto più la probabilità dello stato
proposto è inferiore a quella dello stato corrente.

\begin{osservazione}
Nel caso continuo del nostro algoritmo, $z = \pi(R^*)/\pi(R^{(k-1)}) = e^{\log\rho}$
(con $\log\rho$ definito dalla~\eqref{eq:logrho}). La condizione di accettazione
$\log u < \log\rho$ è equivalente a $u < e^{\log\rho} = z$, ovvero
$u < \min(1, z)$: esattamente la regola $F(z) = \min(1, z)$ derivata qui sopra.
\end{osservazione}

\subsection{Velocità di convergenza}

La velocità con cui la catena converge alla distribuzione stazionaria dipende
dall'\textbf{autocorrelazione della catena}, quantificata dal \emph{tempo di mescolamento}.
Con una proposta uniforme di ampiezza $\delta = 1.5$:

\begin{itemize}
  \item Se $\delta$ è troppo piccola: alta accettazione, ma la catena si muove lentamente
        (esplorazione inefficiente).
  \item Se $\delta$ è troppo grande: bassa accettazione (molte proposte finiscono in zone
        di bassa probabilità), movimenti rari ma grandi.
  \item Il valore ottimale per una target Gaussiana è $\delta \approx 2.38\,\sigma$
        (regola di Gelman et al., 1996), che porta a un tasso di accettazione del 44\%.
\end{itemize}

Con $\delta = 1.5$ e $R(t)$ tipicamente in $[0.5, 3]$, ci si aspetta un tasso di
accettazione nell'intervallo 20--50\% (da verificare empiricamente).

% ===========================================================================
\section{Indipendenza delle stime giornaliere}
% ===========================================================================

Un aspetto fondamentale del Metodo~2 è che la stima di $R(t)$ per ciascun giorno $t$
è \textbf{completamente indipendente} dalle stime degli altri giorni. Per ogni $t$:

\begin{enumerate}
  \item Si fissa $c = C(t)$ (osservazione del giorno $t$).
  \item Si calcola $\lbase(t)$ dalla storia passata $C(t-1), \ldots, C(t-25)$.
  \item Si esegue una catena Metropolis di 5000 passi per campionare $\pi(R \mid c)$.
  \item Si stima $R(t)$ come media della catena post burn-in.
\end{enumerate}

La storia passata entra nella stima \emph{solo} attraverso $\lbase(t)$, che viene
trattata come una costante per l'ottimizzazione del giorno $t$. Questo è molto diverso
da metodi come il filtro di Kalman o il smoothing per serie temporali, in cui la stima
di $R(t)$ dipende anche dall'informazione del futuro.

\textbf{Conseguenza pratica:} il calcolo è parallelizzabile -- si potrebbero eseguire
le $n - \Smax$ catene Metropolis in parallelo su $n - \Smax$ core separati.

% ===========================================================================
\section{Cancellazione della scala}
% ===========================================================================

\subsection{Il problema della sorgente dati}

La metodologia ISS originale utilizza i \emph{casi sintomatici} $S(t)$ (file non
disponibile pubblicamente come CSV scaricabile). Il codice usa invece
\texttt{nuovi\_positivi} DPC, denotato $C_{\mathrm{DPC}}(t)$.

I due dataset differiscono per un fattore di scala (approssimativamente costante):
\begin{equation}
  C_{\mathrm{DPC}}(t) \;\approx\; k \cdot S(t), \qquad k > 1
  \label{eq:scala}
\end{equation}
Il fattore $k > 1$ riflette il fatto che i nuovi positivi includono anche i
pre-sintomatici e gli asintomatici diagnosticati tramite contact tracing.

\subsection{Dimostrazione della cancellazione}

Calcoliamo la base per i dati DPC:
\begin{equation}
  \lbase^{\mathrm{DPC}}(t) = \sum_{s=1}^{25} \phi(s)\, C_{\mathrm{DPC}}(t-s)
  = \sum_{s=1}^{25} \phi(s)\, k\, S(t-s)
  = k \sum_{s=1}^{25} \phi(s)\, S(t-s)
  = k \cdot \lbase^{\mathrm{ISS}}(t)
\end{equation}

Il fattore $k$ moltiplica sia il numeratore che il denominatore dell'MLE:
\begin{equation}
  R_{\mathrm{MLE}}^{\mathrm{DPC}} = \frac{C_{\mathrm{DPC}}(t)}{\lbase^{\mathrm{DPC}}(t)}
  = \frac{k\,S(t)}{k\,\lbase^{\mathrm{ISS}}(t)}
  = \frac{S(t)}{\lbase^{\mathrm{ISS}}(t)}
  = R_{\mathrm{MLE}}^{\mathrm{ISS}}
  \label{eq:cancel}
\end{equation}

Il fattore di scala si cancella esattamente nel rapporto. La stima di $R(t)$ è
identica per le due sorgenti dati (a parità di modello statistico). \qquad $\square$

\subsection{Differenza statistica}

La varianza dell'MLE Poisson è:
\begin{equation}
  \mathrm{Var}[R_{\mathrm{MLE}}] \approx \frac{R(t)}{\lbase(t)} = \frac{R(t)^2}{C(t)}
\end{equation}
(approssimazione di Cramér--Rao). Poiché $C_{\mathrm{DPC}} = k\,S$, i dati DPC producono
$\lbase^{\mathrm{DPC}} = k\,\lbase^{\mathrm{ISS}}$ e quindi:
\begin{equation}
  \mathrm{Var}[R^{\mathrm{DPC}}_{\mathrm{MLE}}] = \frac{R(t)}{k\,\lbase^{\mathrm{ISS}}(t)}
  = \frac{1}{k}\, \mathrm{Var}[R^{\mathrm{ISS}}_{\mathrm{MLE}}]
\end{equation}
Usare i dati DPC (più numerosi per il fattore $k > 1$) produce
\textbf{intervalli di confidenza più stretti}: un vantaggio inatteso della sostituzione.

% ===========================================================================
\section{Post-processing: media mobile causale}
% ===========================================================================

La stima $\widehat{R}(t)$ grezza presenta ancora oscillazioni per due motivi:

\begin{enumerate}
  \item \textbf{Rumore osservativo}: le segnalazioni epidemiologiche hanno forti
        oscillazioni settimanali (meno tamponi nel weekend, aggiornamenti batch il lunedì).
        Questo rumore entra in $C(t)$ e quindi in $\widehat{R}(t)$.

  \item \textbf{Fluttuazioni Monte Carlo}: la catena Metropolis di 4500 passi ha una
        varianza residua piccola ($\sigma \approx 3 \times 10^{-3}$, trascurabile rispetto
        al rumore osservativo).
\end{enumerate}

Si applica una \textbf{media mobile causale a 7 giorni}:
\begin{equation}
  R_{\mathrm{MC}}(t) \;=\; \frac{1}{7}\sum_{k=0}^{6} \widehat{R}(t-k)
  \;=\; \frac{\widehat{R}(t) + \widehat{R}(t-1) + \cdots + \widehat{R}(t-6)}{7}
  \label{eq:smooth}
\end{equation}

Il termine \textbf{causale} (o \emph{one-sided}) significa che la finestra guarda
solo al passato e al presente: non vengono mai usati valori futuri $\widehat{R}(t+1), \ldots$.
In un sistema di sorveglianza in tempo reale, i dati del futuro non sono disponibili;
la scelta causale è quindi l'unica fisicamente sensata.

Il periodo di 7 giorni è scelto per attenuare il ciclo settimanale delle segnalazioni
(periodo di 7 giorni): la media mobile elimina esattamente le componenti periodiche
il cui periodo divide la finestra.

% ===========================================================================
\section{Il ritardo strutturale: confronto con il dato INFN}
% ===========================================================================

I dati ufficiali INFN/CovidStat usano la \textbf{data di inizio sintomi} (stimata tramite
un modello di nowcasting). Il Metodo~2 usa la data di referto DPC, sistematicamente
più tardiva di circa 5--10 giorni.

Paradossalmente, il Metodo~2 produce una stima di $R(t)$ che \textbf{anticipa} il dato
INFN di circa una settimana. Il motivo è sottile:

\begin{itemize}
  \item Il dato INFN al giorno $t$ include anche le \emph{correzioni retroattive} dei
        giorni precedenti (i casi con data di inizio sintomi nel passato ma segnalati
        con ritardo). Questo introduce un ritardo di pubblicazione di circa una settimana.
  \item Il Metodo~2 usa i casi \emph{così come appaiono nel database DPC} al giorno $t$,
        senza aspettare correzioni retroattive. La stima è quindi più tempestiva e, poiché
        $\lbase(t)$ dipende solo dai casi passati rispetto a $t$, non subisce aggiornamenti
        retroattivi.
\end{itemize}

% ===========================================================================
\section{Confronto sistematico con il Metodo 1 (SIR)}
% ===========================================================================

\begin{center}
\renewcommand{\arraystretch}{1.6}
\begin{tabular}{p{4cm}p{5cm}p{5cm}}
\toprule
\textbf{Aspetto} & \textbf{Metodo 1 --- SIR} & \textbf{Metodo 2 --- Metropolis} \\
\midrule
Modello & Deterministico (formula approssimata) & Probabilistico (Poisson) \\
Dipendenza da $\gamma$ & Esplicita: $R \approx 1 + T_R\,\Delta I/I$ & Assente \\
Dist. tempi generazione & Implicita (finestra fissa $T_R$) & Esplicita: Gamma($\alpha$=1.87, $\theta$=3.57) \\
Ritardo vs INFN & $+1$ settimana (data referto $>$ data sintomi) & $-1$ settimana (anticipo) \\
Errore max vs INFN & $0.2$--$0.6$ & $0.20$--$0.27$ \\
Sensibilità al rumore & Alta (senza pre-smoothing) & Moderata (distribuita su 25 giorni) \\
Costo computazionale & Basso (formula algebrica) & Alto ($\sim$5 min su MATLAB) \\
Funziona senza $T_R$ & No & Sì \\
\bottomrule
\end{tabular}
\end{center}

% ===========================================================================
\section{Riepilogo completo}
% ===========================================================================

\begin{riepilogo}[title=Pipeline del Metodo 2]
\begin{align*}
&\textbf{Input:} \quad C(1), C(2), \ldots, C(n) \quad \text{[nuovi\_positivi DPC]} \\
&\quad \phi(s) \quad s=1,\ldots,25 \quad \text{[Gamma(1.87, 3.57), discretizzata, normalizzata]} \\[6pt]
&\textbf{Per ogni } t = 26, 27, \ldots, n\textbf{:} \\
&\quad \lbase(t) = \sum_{s=1}^{25} \phi(s)\, C(t-s) \\
&\quad \text{Catena Metropolis, } N_{\mathrm{iter}}=5000,\; N_{\mathrm{burnin}}=500,\; \delta=1.5\text{:} \\
&\qquad \text{target: } \pi(R) \propto \exp\!\bigl[C(t)\log R - R\lbase(t)\bigr] \\
&\quad \widehat{R}(t) = \frac{1}{4500}\sum_{k=501}^{5000} R^{(k)} \\[6pt]
&\textbf{Post-processing:} \\
&\quad R_{\mathrm{MC}}(t) = \frac{1}{7}\sum_{k=0}^{6} \widehat{R}(t-k) \quad \text{[media mobile causale 7 giorni]}
\end{align*}
\end{riepilogo}

\end{document}
